\chapter{Carnets d'ordres on-chain et DEX de perpétuels}\label{ch:m3:on-chain-order-books-and-perpetual-dexs}

Sur une \omterm{def:m3:automated-market-makers:amm}{plateforme décentralisée}, chaque ordre, chaque annulation et chaque liquidation d'un marché de \omterm{def:m3:perpetual-futures:perp}{contrats perpétuels}
est une transaction incluse dans un bloc. Les \omterm{def:m3:blockchains-for-traders:pow}{validateurs} de la plateforme s'accordent non seulement sur ce qui s'est
produit, mais aussi sur l'ordre dans lequel cela s'est produit, et cet ordre détermine qui est exécuté~: un teneur de marché
dont la cotation n'est plus à jour envoie une annulation, un arbitragiste envoie un ordre pour frapper la cotation, et les
deux arrivent dans le même bloc. Sur une \omterm{def:m3:centralised-exchanges:cex}{plateforme centralisée}, c'est l'horloge du moteur d'appariement qui tranche la
course. Sur une chaîne, c'est une règle inscrite dans le protocole. Ce chapitre explique pourquoi il est difficile de placer
un carnet d'ordres sur une chaîne, examine en détail la conception de la plateforme qui est allée le plus loin dans cette
voie, analyse ce qui a mal tourné lorsque les marchés ont été soumis à des tensions, et compare les autres modèles~: les
carnets d'ordres tenus hors chaîne par des \omterm{def:m3:blockchains-for-traders:pow}{validateurs}, et les \omterm{def:m3:on-chain-order-books-and-perpetual-dexs:oracle}{plateformes à prix d'oracle}, qui renoncent totalement aux carnets d'ordres.

\section{Pourquoi un carnet d'ordres on-chain est difficile}

\begin{definition}[DEX de perpétuels, carnet d'ordres on-chain]\label{def:m3:on-chain-order-books-and-perpetual-dexs:perpdex}
Un \emph{DEX de perpétuels}\index{DEX de perpétuels} est une plateforme d'échange décentralisée pour les contrats à terme
perpétuels, dont la marge, les positions et le règlement-livraison sont détenus et calculés par l'état d'une \omterm{def:m3:blockchains-for-traders:chain}{blockchain}
plutôt que par une société. Un \emph{carnet d'ordres on-chain}\index{carnet d'ordres on-chain} est un carnet d'ordres à
cours limité dont les ordres, annulations et appariements sont des transactions exécutées et enregistrées par le consensus de la chaîne.
\end{definition}

Un teneur de marché sur une \omterm{def:m3:centralised-exchanges:cex}{plateforme centralisée} envoie des milliers de messages par seconde et annule la plupart de ses
ordres~; un message ne coûte rien et est traité en quelques microsecondes. Sur une chaîne à usage général, chaque message
est une transaction qui paie du \omterm{def:m3:blockchains-for-traders:gas}{gas} (\cref{ch:m3:blockchains-for-traders}) et attend un bloc, distant de quelques secondes~:
coter coûte de l'argent, les annulations sont lentes, et quiconque surveille le \omterm{def:m3:blockchains-for-traders:mempool}{mempool} voit venir l'annulation. Trois
conséquences en découlent. La chaîne doit être construite pour cet usage, avec des blocs produits en fractions de seconde
et aucun \omterm{def:m3:blockchains-for-traders:gas}{gas} par message pour les actions de trading. L'ordre des actions à l'intérieur d'un bloc doit être fixé par le
protocole, car le proposant du bloc pourrait sinon vendre le droit de passer en premier. Et toutes les fonctions de
l'échange qu'une société assumerait (l'indice, le \omterm{def:m3:perpetual-futures:index}{prix de marquage}, la liquidation, le filet de sécurité) doivent être
assurées par le protocole et ses \omterm{def:m3:blockchains-for-traders:pow}{validateurs}, de manière publique.

\begin{definition}[App-chain]\label{def:m3:on-chain-order-books-and-perpetual-dexs:appchain}
Une \emph{app-chain}\index{app-chain} est une \omterm{def:m3:blockchains-for-traders:chain}{blockchain} conçue pour faire fonctionner une seule application, ici une
plateforme d'échange, avec ses propres \omterm{def:m3:blockchains-for-traders:pow}{validateurs}, ses paramètres de consensus et ses règles de transaction choisis pour
cette application plutôt que pour le calcul généraliste.
\end{definition}

\section{Hyperliquid en détail}

Hyperliquid est une \omterm{def:m3:on-chain-order-books-and-perpetual-dexs:appchain}{app-chain} dont l'état comprend les carnets d'ordres, la marge et les positions de ses marchés de
perpétuels et au comptant (elle appelle cette partie HyperCore), aux côtés d'un environnement de \omterm{def:m3:blockchains-for-traders:contract}{smart contracts} à usage général.

\begin{dated}{2026-09}{Consensus, latence et débit}\label{dat:m3:on-chain-order-books-and-perpetual-dexs:hl}
La documentation d'Hyperliquid décrit une chaîne en \omterm{def:m3:blockchains-for-traders:pow}{preuve d'enjeu} (proof of stake), sécurisée par HyperBFT, une variante du consensus
HotStuff, avec les carnets d'ordres et l'état de la marge conservés dans l'état de la chaîne et non hors chaîne. Elle annonce une latence médiane
de bout en bout de 0.2 seconde et un 99e percentile de 0.9 seconde pour un ordre émis depuis un client co-localisé, ainsi qu'un
débit d'environ 200\,000 ordres par seconde, limité par l'exécution. Elle précise qu'il n'existe aucun programme de teneurs
de marché désignés (DMM), ni remises particulières, ni frais spécifiques, ni avantages de latence.
\end{dated}

\begin{definition}[Règle de priorité dans le bloc]\label{def:m3:on-chain-order-books-and-perpetual-dexs:priority}
Une \emph{règle de priorité dans le bloc}\index{règle de priorité dans le bloc} est la règle du protocole qui trie les actions de trading
d'un bloc avant de les appliquer aux carnets d'ordres, en passant outre l'ordre dans lequel le proposant du bloc les a reçues
ou inscrites.
\end{definition}

\begin{dated}{2026-09}{L'ordonnancement dans le bloc d'une plateforme}\label{dat:m3:on-chain-order-books-and-perpetual-dexs:order}
La documentation d'Hyperliquid indique qu'à l'intérieur d'un bloc, les actions sont triées dans chaque lot de consensus (en
général un ou deux par bloc) comme suit~: d'abord les actions qui n'envoient aucun ordre GTC ou IOC à un carnet~; ensuite les annulations~; enfin les actions qui envoient au
moins un ordre GTC ou IOC. Au sein de chaque catégorie, les actions conservent l'ordre proposé par le proposant du bloc. Ses carnets
sont appariés selon la priorité prix-temps, avec des prix sur une grille de tick et des tailles sur une grille de lot~; les \omterm{def:m3:centralised-exchanges:postonly}{ordres post-only} (ALO) restent au carnet sans s'exécuter.
\end{dated}

\begin{omfigure}
\centering
\begin{tikzpicture}[font=\small,
  c/.style={draw,thick,rounded corners=2pt,align=center,minimum height=0.9cm,text width=3.3cm},
  arr/.style={-{Stealth[length=2.2mm]},thick}]
\node[draw,thick,rounded corners=2pt,align=left,text width=4.3cm,fill=gray!6] (in) at (-0.4,0) {\footnotesize actions telles que proposées~:\\
achat IOC (preneur)\\annulation d'un vendeur périmé (teneur)\\nouvel ordre vendeur post-only (teneur)\\vente GTC (trader)};
\node[c,draw=omDef,fill=omDef!8] (c1) at (5.6,1.3) {1.\ sans GTC/IOC~:\\cotations post-only};
\node[c,draw=omThm,fill=omThm!8] (c2) at (5.6,0) {2.\ annulations};
\node[c,draw=omExo,fill=omExo!8] (c3) at (5.6,-1.3) {3.\ ordres GTC et IOC};
\node[c,draw=omProp,fill=omProp!8] (bk) at (10.6,0) {carnet d'ordres,\\priorité prix-temps};
\draw[arr] (in.east) -- (c1.west); \draw[arr] (in.east) -- (c2.west); \draw[arr] (in.east) -- (c3.west);
\draw[arr] (c1.east) -- (bk.west); \draw[arr] (c2.east) -- (bk.west); \draw[arr] (c3.east) -- (bk.west);
\end{tikzpicture}

\omcaption{Une \omterm{def:m3:on-chain-order-books-and-perpetual-dexs:priority}{règle de priorité dans le bloc} du type décrit par l'\cref{dat:m3:on-chain-order-books-and-perpetual-dexs:order}~: les
actions du bloc sont réparties en trois classes et appliquées classe par classe, dans l'ordre proposé au sein de chaque classe. L'annulation d'un teneur
et la nouvelle cotation post-only sont appliquées avant tout ordre d'un preneur dans le même bloc. Schéma.}\label{fig:m3:on-chain-order-books-and-perpetual-dexs:classes}
\end{omfigure}

La règle constitue un transfert délibéré de valeur des preneurs vers les teneurs. Sur une \omterm{def:m3:centralised-exchanges:cex}{plateforme centralisée}, un preneur qui réagit plus vite que
l'annulation d'un teneur capte la cotation périmée~; ici, si les deux atteignent le même bloc, l'annulation l'emporte, de même que la nouvelle
cotation post-only du teneur. Les teneurs peuvent coter plus serré parce qu'ils sont moins souvent captés, et un preneur qui veut frapper une cotation périmée doit
entrer dans un bloc antérieur à l'annulation du teneur. Le tutoriel rejoue mille blocs~: lorsque le preneur arrive en premier
la moitié du temps, l'ordre d'arrivée lui permet de capter 610 lots, soit 940 ticks de la valeur du teneur~; avec la règle par classe, aucun.

\begin{omfigure}
\begin{tikzpicture}
\begin{axis}[width=11cm,height=5cm,
  xlabel={bloc},ylabel={perte cumulée du teneur, ticks},
  tick label style={font=\small},label style={font=\small},
  xmin=0,xmax=1000,ymin=0,ymax=1000,grid=major,grid style={gray!25},
  legend columns=2,legend style={font=\footnotesize,at={(0.5,-0.3)},anchor=north,draw=none,fill=none,/tikz/every even column/.append style={column sep=8pt}},
  table/col sep=comma]
\addplot[omThm,very thick] table[x=block,y=arrival]{figdata/markets-3/21-on-chain-order-books-and-perpetual-dexs/pickoff.csv};\addlegendentry{ordre d'arrivée}
\addplot[omDef,very thick,dashed] table[x=block,y=cancels_first]{figdata/markets-3/21-on-chain-order-books-and-perpetual-dexs/pickoff.csv};\addlegendentry{cotations et annulations avant les preneurs}
\end{axis}
\end{tikzpicture}

\omcaption{Pertes d'un teneur de marché liées aux captures de cotations périmées sur mille blocs selon deux règles de priorité,
la valeur de référence évoluant par variations normales de deux ticks par bloc et l'ordre du preneur arrivant avant l'annulation du teneur la moitié du temps. Avec la règle par
classe, l'annulation et la nouvelle cotation du teneur passent toujours en premier. Réplication synthétique. Données~: le tutoriel du chapitre.}\label{fig:m3:on-chain-order-books-and-perpetual-dexs:pickoff}
\end{omfigure}

Le reste de l'échange est lui aussi on-chain. Les prix proviennent des \omterm{def:m3:blockchains-for-traders:pow}{validateurs} plutôt que des seules transactions de la plateforme, afin que les
carnets peu fournis d'une jeune plateforme ne puissent servir à déplacer ses propres références.

\begin{dated}{2026-09}{Oracle, prix de marquage et liquidation sur une même plateforme}\label{dat:m3:on-chain-order-books-and-perpetual-dexs:oracle}
Les \omterm{def:m3:blockchains-for-traders:pow}{validateurs} d'Hyperliquid publient chacun, environ toutes les trois secondes, un prix d'\omterm{def:m3:blockchains-for-traders:contract}{oracle} au comptant par \omterm{def:m3:perpetual-futures:perp}{contrat perpétuel}~: une médiane pondérée des prix milieu
au comptant sur Binance, OKX, Bybit, Kraken, KuCoin, Gate.io, MEXC et Hyperliquid, avec des poids de 3, 2, 2, 1, 1, 1, 1 et 1 (sources externes
seulement pour les actifs dont la liquidité au comptant principale se trouve ailleurs)~; la chambre de compensation utilise la médiane des prix des \omterm{def:m3:blockchains-for-traders:pow}{validateurs} pondérée par l'enjeu.
Le \omterm{def:m3:perpetual-futures:index}{prix de marquage} est la médiane du prix d'\omterm{def:m3:blockchains-for-traders:contract}{oracle} majoré d'une moyenne mobile de 150~secondes de la base, de la médiane du meilleur cours acheteur
du meilleur cours vendeur et de la dernière transaction de la plateforme, et d'une médiane pondérée des prix milieu des \omterm{def:m3:perpetual-futures:perp}{contrats perpétuels} sur cinq plateformes externes. La marge de maintenance vaut la moitié de la marge initiale
au levier maximal (1.25~\% à 16.7~\%)~; les liquidations sont d'abord envoyées au carnet d'ordres sous forme d'ordres au marché (20~\% à la fois pour les positions
supérieures à 100\,000 USDC)~; sous les deux tiers de la marge de maintenance, une position est reprise par le coffre de liquidation, une stratégie du
coffre du protocole nommée HLP, dont les déposants partagent le P\&L.
\end{dated}

\begin{definition}[Coffre de liquidité]\label{def:m3:on-chain-order-books-and-perpetual-dexs:vault}
Un \emph{coffre de liquidité}\index{coffre de liquidité} est un ensemble de capitaux apportés par des déposants et géré par un protocole, ou par un gestionnaire selon les règles du protocole,
qui fait office de teneur de marché, reprend les positions liquidées ou apporte du capital à une plateforme, et en partage le P\&L entre ses déposants.
\end{definition}

Un coffre qui sert de filet aux liquidations est le \omterm{def:m3:margin-liquidation-and-loss-allocation:fund}{fonds d'assurance} de la plateforme et son teneur de marché de dernier recours, dont la propriété revient à quiconque y dépose des fonds.
En moyenne, ses déposants encaissent la marge de sécurité des liquidations et supportent les pertes extrêmes lorsqu'une position qu'il reprend ne peut être clôturée à un prix raisonnable.

\section{Incidents}

Les tests de résistance éprouvent la conception là où le personnel d'une \omterm{def:m3:centralised-exchanges:cex}{plateforme centralisée} interviendrait. Lors de la vente massive du 10 octobre 2025
(\cref{ch:m3:margin-liquidation-and-loss-allocation}), CoinDesk Research a mesuré la plus forte contraction des positions ouvertes sur Hyperliquid,
passant de 14~milliards de dollars à 6~milliards de dollars, soit une baisse de 57~\%. Le mécanisme d'\omterm{def:m3:margin-liquidation-and-loss-allocation:adl}{auto-deleveraging} (ADL) de la plateforme a clôturé des positions gagnantes d'autres traders, et une analyse
menée par Tarun Chitra de Gauntlet, rapportée par Protos, affirmait que son classement fondé sur la file d'attente avait clôturé environ 650~millions de dollars de positions de plus que ce qui était
nécessaire. Plus tôt, le 12 mars 2025, un trader détenant une position longue à effet de levier de 50 fois sur 113\,000 ether a retiré de la marge jusqu'à ce que la position soit
liquidée dans le coffre~: le trader a conservé environ 1.8~million de dollars de bénéfice et le coffre a perdu environ 4~millions de dollars~; la plateforme a alors ramené
son levier maximal à 40 fois sur le bitcoin et à 25 fois sur l'ether.

Un deuxième type d'incident est propre à une plateforme dont le filet de sécurité est un coffre public. Un trader ouvre une position importante sur un marché peu liquide et la laisse
se faire liquider pour que le coffre la reprenne, puis fait bouger les prix au comptant externes qui alimentent l'\omterm{def:m3:blockchains-for-traders:contract}{oracle}, sur de plus petites plateformes où cela est
peu coûteux~; le coffre, qui détient une position qu'il ne peut clôturer sans déplacer le carnet, valorise une perte croissante. Les défenses du protocole sont un
\omterm{def:m3:on-chain-order-books-and-perpetual-dexs:oicap}{plafond de positions ouvertes} par actif et par compte, des paliers de marge qui relèvent la marge de maintenance avec la taille, et, en dernier recours, la gouvernance~: les
\omterm{def:m3:blockchains-for-traders:pow}{validateurs} peuvent voter pour retirer un marché de la cote.

\begin{definition}[Plafond de positions ouvertes]\label{def:m3:on-chain-order-books-and-perpetual-dexs:oicap}
Un \emph{plafond de positions ouvertes}\index{plafond de positions ouvertes} est une limite sur les positions ouvertes totales d'un marché, ou sur la position d'un seul compte dans celui-ci,
au-delà de laquelle la plateforme rejette les ordres qui augmenteraient ces positions ouvertes.
\end{definition}

\begin{proposition}[La perte d'un coffre qui reprend une position soumise à un squeeze]\label{prop:m3:on-chain-order-books-and-perpetual-dexs:squeeze}
Un coffre reprend une position vendeuse de notionnel $N$ dont une fraction $b$ reste en marge, et le prix monte ensuite d'une fraction $s$ avant que le
coffre puisse clôturer. Sa perte vaut $N(s - b)^+$. Pour que la perte reste inférieure à $L$ pour des squeezes allant jusqu'à $s_{\max}$, les positions qu'il peut reprendre doivent être plafonnées à
$N \le L/(s_{\max} - b)$.
\end{proposition}

\begin{proof}
La position vendeuse perd $Ns$ et la marge reprise couvre $Nb$ de cette perte~; la borne s'obtient en résolvant $N(s_{\max} - b) \le L$.
\end{proof}

\begin{omfigure}
\begin{tikzpicture}
\begin{axis}[width=11cm,height=5cm,
  xlabel={squeeze après la reprise, \%},ylabel={perte du coffre, millions de dollars},
  tick label style={font=\small},label style={font=\small},
  xmin=0,xmax=500,ymin=0,ymax=50,grid=major,grid style={gray!25},
  legend columns=3,legend style={font=\footnotesize,at={(0.5,-0.3)},anchor=north,draw=none,fill=none,/tikz/every even column/.append style={column sep=8pt}},
  table/col sep=comma]
\addplot[omThm,very thick] table[x=squeeze_pct,y=nocap]{figdata/markets-3/21-on-chain-order-books-and-perpetual-dexs/squeeze.csv};\addlegendentry{10~millions de dollars hérités}
\addplot[omDef,very thick,dashed] table[x=squeeze_pct,y=cap2]{figdata/markets-3/21-on-chain-order-books-and-perpetual-dexs/squeeze.csv};\addlegendentry{plafond de 2~millions de dollars}
\addplot[omExo,very thick,dotted] table[x=squeeze_pct,y=cap1]{figdata/markets-3/21-on-chain-order-books-and-perpetual-dexs/squeeze.csv};\addlegendentry{plafond de 1~million de dollars}
\end{axis}
\end{tikzpicture}

\omcaption{Perte d'un coffre qui reprend une position vendeuse sur un marché dont le levier maximal est de 3 fois (marge de maintenance de 16.7~\%, reprise aux deux tiers de celle-ci,
laissant 11.1~\% du notionnel en marge), en fonction de la hausse qui suit (\cref{prop:m3:on-chain-order-books-and-perpetual-dexs:squeeze}). Positions
illustratives. Données~: le tutoriel du chapitre.}\label{fig:m3:on-chain-order-books-and-perpetual-dexs:squeeze}
\end{omfigure}

\begin{dated}{2026-09}{Retrait de la cote par les validateurs}\label{dat:m3:on-chain-order-books-and-perpetual-dexs:delist}
La documentation d'Hyperliquid indique que les \omterm{def:m3:blockchains-for-traders:pow}{validateurs} votent sur le retrait de la cote des perpétuels gérés par des \omterm{def:m3:blockchains-for-traders:pow}{validateurs}~; si c'est le cas, les perpétuels se règlent au
prix d'\omterm{def:m3:blockchains-for-traders:contract}{oracle} au comptant pondéré dans le temps sur une heure, avant l'heure prévue du vote~; toutes les positions sont réglées et les ordres en attente annulés.
Le scénario décrit plus haut s'est produit le 26 mars 2025~: une position vendeuse sur JELLY est passée au coffre, HLP, par liquidation, puis des achats au comptant ailleurs ont porté la perte latente du coffre
à 13.5~millions de dollars, et les \omterm{def:m3:blockchains-for-traders:pow}{validateurs} ont voté le retrait de la cote du marché, qui a été réglé à 0.0095 alors qu'un prix d'environ 0.50 était encore transmis à
l'\omterm{def:m3:blockchains-for-traders:contract}{oracle}.
\end{dated}

Un vote qui règle un marché à un prix est l'équivalent décentralisé des pouvoirs d'urgence d'une bourse. Il peut protéger un coffre contre un squeeze~; il
montre aussi que les règles du marché peuvent être modifiées par les \omterm{def:m3:blockchains-for-traders:pow}{validateurs} alors que des positions sont ouvertes, un risque de gouvernance qu'un trader intègre dans ses prix de la même manière
que la discrétion d'une \omterm{def:m3:centralised-exchanges:cex}{plateforme centralisée}.

\section{Conceptions app-chain et à prix d'oracle}

Deux autres conceptions échangent certaines propriétés du \omterm{def:m3:on-chain-order-books-and-perpetual-dexs:perpdex}{carnet d'ordres on-chain} contre de la performance ou de la simplicité.

\begin{dated}{2026-09}{Une app-chain avec un carnet hors chaîne}\label{dat:m3:on-chain-order-books-and-perpetual-dexs:dydx}
La documentation de dYdX décrit sa chaîne («~v4~») comme une chaîne à \omterm{def:m3:blockchains-for-traders:pow}{preuve d'enjeu} bâtie sur CometBFT et le Cosmos SDK, dont les \omterm{def:m3:blockchains-for-traders:pow}{validateurs} conservent le carnet d'ordres
en mémoire, hors chaîne et hors du consensus~; lorsque des ordres se croisent, le proposant du bloc ajoute les appariements à son bloc proposé, et un bloc est validé lorsque
deux tiers de l'enjeu l'approuvent. Les proposants se relaient selon un tour de rôle pondéré par l'enjeu.
\end{dated}

\begin{definition}[Plateforme à prix d'oracle]\label{def:m3:on-chain-order-books-and-perpetual-dexs:oracle}
Une \emph{plateforme à prix d'oracle}\index{plateforme à prix d'oracle} est un \omterm{def:m3:on-chain-order-books-and-perpetual-dexs:perpdex}{DEX de perpétuels} ou au comptant qui n'a pas de carnet d'ordres~: les traders ouvrent et clôturent des positions
contre un \omterm{def:m3:automated-market-makers:amm}{pool de liquidité} à un prix fourni par un \omterm{def:m3:blockchains-for-traders:contract}{oracle}, avec des frais et des coûts d'impact sur le prix fixés par le protocole, et les déposants du pool prennent la
contrepartie de chaque transaction.
\end{definition}

\begin{dated}{2026-09}{Une plateforme de perpétuels à prix d'oracle}\label{dat:m3:on-chain-order-books-and-perpetual-dexs:gmx}
La documentation de GMX décrit une \omterm{def:m3:automated-market-makers:amm}{plateforme décentralisée} au comptant et de perpétuels sur Arbitrum, Avalanche et MegaETH, avec un effet de levier pouvant atteindre 100 fois, qui achemine
chaque ordre contre ses \omterm{def:m3:automated-market-makers:amm}{pools de liquidité} et cote au \omterm{def:m3:perpetual-futures:index}{prix de l'indice} d'\omterm{def:m3:blockchains-for-traders:contract}{oracle} plutôt que de s'appuyer sur un carnet d'ordres ou des teneurs de marché externes~; il utilise
les \omterm{def:m3:blockchains-for-traders:contract}{oracles} Chainlink Data Streams, et les fournisseurs de liquidité perçoivent 63~\% des frais sur les transactions, les liquidations, les emprunts et les swaps sur Arbitrum et Avalanche.
\end{dated}

\begin{omfigure}
\centering
\begin{tikzpicture}[font=\footnotesize,
  h/.style={draw,thick,rounded corners=2pt,align=center,text width=2.9cm,minimum height=0.8cm},
  r/.style={align=center,text width=2.9cm}]
\node[h,draw=omDef,fill=omDef!8] at (0,0) {plateforme\\centralisée};
\node[h,draw=omThm,fill=omThm!8] at (3.5,0) {carnet d'ordres\\on-chain};
\node[h,draw=omProp,fill=omProp!8] at (7,0) {carnet hors chaîne,\\appariements on-chain};
\node[h,draw=omExo,fill=omExo!8] at (10.5,0) {pool à prix\\d'oracle};
\node[r,anchor=north] at (0,-0.6) {carnet~: société\\ordre~: son horloge\\conservation~: société};
\node[r,anchor=north] at (3.5,-0.6) {carnet~: chaîne\\ordre~: règle\\conservation~: chaîne};
\node[r,anchor=north] at (7,-0.6) {carnet~: validateurs\\ordre~: proposant\\conservation~: chaîne};
\node[r,anchor=north] at (10.5,-0.6) {carnet~: aucun\\prix~: oracle\\conservation~: chaîne};
\end{tikzpicture}

Quatre conceptions d'une plateforme de produits dérivés, selon qui détient le carnet, qui décide de l'ordre des actions et qui détient le collatéral. Schématique.\label{fig:m3:on-chain-order-books-and-perpetual-dexs:designs}
\end{omfigure}

La conception à prix d'\omterm{def:m3:blockchains-for-traders:contract}{oracle} supprime les jeux de carnet d'ordres et les remplace par des jeux d'\omterm{def:m3:blockchains-for-traders:contract}{oracle}~: le pool perd chaque fois qu'un trader connaît le prix mieux que l'\omterm{def:m3:blockchains-for-traders:contract}{oracle}, comme
un \omterm{def:m3:automated-market-makers:amm}{teneur de marché automatisé} perd face aux arbitragistes (\cref{ch:m3:automated-market-makers}), et il doit se défendre par des frais, des délais et des limites sur
les positions ouvertes. Le carnet hors chaîne restaure la rapidité, mais confie l'ordonnancement des appariements au proposant de chaque bloc. Pour une société de trading, les questions sont celles de
toute plateforme (\cref{ch:m3:centralised-exchanges}), avec d'autres réponses~: qui ordonne ses messages, ce qu'il advient de son collatéral si la plateforme tombe en panne, et qui
peut modifier les règles pendant qu'elle est positionnée.

\section{Tutoriel~: deux ordonnancements d'un même bloc}\label{tut:m3:on-chain-order-books-and-perpetual-dexs:replay}

\begin{tutorial}
\textbf{Objectif.} Rejouer des blocs de cotations de teneur de marché, d'annulations et d'ordres preneurs à travers un carnet à priorité prix-temps, selon l'ordre d'arrivée puis selon la règle de classe, et
compter les cotations périmées capturées. \textbf{État final~:} \cref{fig:m3:on-chain-order-books-and-perpetual-dexs:pickoff,fig:m3:on-chain-order-books-and-perpetual-dexs:squeeze}.
\begin{tutsteps}
\item \textbf{La règle.} Trier les actions d'un bloc par classe, en conservant l'ordre d'arrivée au sein de chaque classe.
\omcode{code/firm/blockbook/firm_blockbook.py}{83}{96}{Ordonnancement dans le bloc et exécution du bloc.}\label{lst:m3:on-chain-order-books-and-perpetual-dexs:order}
\item \textbf{Le rejeu.} Un teneur de marché recote à chaque bloc~; un preneur attaque les cotations que la fair value a dépassées.
\omcode{code/markets-3/21-on-chain-order-books-and-perpetual-dexs/python/m3_perpdex.py}{12}{41}{Rejeu des captures sur cotations périmées.}\label{lst:m3:on-chain-order-books-and-perpetual-dexs:replay}
\item \textbf{Exécuter} \texttt{replay} avec les règles \texttt{arrival} et \texttt{cancels\_first}, puis \texttt{fig\_perpdex.py}.
\end{tutsteps}
\textbf{Pour aller plus loin.} Faire arriver le preneur en premier dans 90~\% des cas~; laisser l'annulation du teneur de marché tomber dans le bloc suivant avec une certaine probabilité et observer ce que la
règle de classe protège encore~; ajouter un second teneur de marché qui n'annule jamais.
\end{tutorial}

\section{Construction~: le carnet groupé par bloc}\label{bld:m3:on-chain-order-books-and-perpetual-dexs:blockbook}

\begin{build}
\bfield{Rôle} Pour négocier sur des plateformes à blocs, la petite société doit prédire la façon dont ses actions seront ordonnancées et appariées~; pour faire du backtest, elle doit rejouer
les blocs à l'identique~; et son risque doit intégrer les \omterm{def:m3:on-chain-order-books-and-perpetual-dexs:oicap}{plafonds de positions ouvertes} de la plateforme ainsi que son \omterm{def:m3:perpetual-futures:index}{prix de marquage} (mark price).
\bfield{Interface} \texttt{Action(kind, trader, side, price, qty, oid)} avec les types \texttt{alo}, \texttt{gtc}, \texttt{ioc}, \texttt{cancel}~;
\texttt{Book(oi\_cap)} avec \texttt{apply(action)} et \texttt{best()}~; \texttt{order\_block(actions, rule)} pour \texttt{arrival} et
\texttt{cancels\_first}~; \texttt{run\_block(book, actions, rule)}~; \texttt{mark\_price(oracle\_basis, book\_median, external)}, la médiane des trois. Utilise \texttt{firm.perp} et
\texttt{firm.liquidation} pour la marge.
\bfield{Règles} Ticks et lots entiers~; tri stable, de sorte que les rejeux soient déterministes~; les \omterm{def:m3:centralised-exchanges:postonly}{ordres post-only} qui croiseraient sont rejetés~; les ordres qui porteraient une
position au-delà du plafond sont rejetés.
\bfield{Tests d'acceptation} \texttt{code/firm/blockbook/tests/}~: priorité prix-temps et rejet des post-only~; une cotation périmée capturée selon l'ordre d'arrivée et protégée
selon la règle de classe~; le \omterm{def:m3:on-chain-order-books-and-perpetual-dexs:oicap}{plafond de positions ouvertes}~; la médiane du mark.
\bfield{Pour aller plus loin} Consensus par lots à l'intérieur d'un bloc~; contrôles de marge à la saisie de l'ordre et à l'exécution du côté au repos~; ordres de liquidation entrant dans le bloc sous forme d'IOC.
\end{build}

\begin{omsources}
\item Documentation Hyperliquid~: HyperCore overview, order book, oracle, robust price indices, liquidations, protocol vaults, delisting, market making.
Consultée en septembre 2026.
\item Documentation dYdX, «~Intro to dYdX Chain Architecture~»~; documentation GMX, introduction. Consultées en septembre 2026.
\item CoinDesk Research, 17 octobre 2025~; Protos, 10 décembre 2025 (voir le \cref{ch:m3:margin-liquidation-and-loss-allocation}).
\item CoinDesk, «~Hyperliquid Loses \$4M After Whale's Over \$200M Ether Trade Unwinds~», 12 mars 2025~; «~HyperLiquid Delists JELLY
After Vault Squeezed in \$13M Tussle~», 26 mars 2025.
\end{omsources}

\section{Exercices}

\begin{exercise}[$\star$]\label{exo:m3:on-chain-order-books-and-perpetual-dexs:1}
Dans un même bloc arrivent, dans cet ordre~: un IOC à l'achat à 101, une annulation du cours vendeur à 101 et un post-only à la vente à 103. Qu'est-ce qui s'exécute selon l'ordre d'arrivée, et selon la
règle de classe~?
\end{exercise}

\begin{exercise}[$\star$]\label{exo:m3:on-chain-order-books-and-perpetual-dexs:2}
Pourquoi une chaîne généraliste à blocs de 12 secondes et à \omterm{def:m3:blockchains-for-traders:gas}{gas} par transaction ne peut-elle pas accueillir un carnet d'ordres compétitif~?
\end{exercise}

\begin{exercise}[$\star$]\label{exo:m3:on-chain-order-books-and-perpetual-dexs:3}
Contre quoi protège la médiane pondérée par l'enjeu des prix d'\omterm{def:m3:blockchains-for-traders:contract}{oracle} des \omterm{def:m3:blockchains-for-traders:pow}{validateurs}~?
\end{exercise}

\begin{exercise}[$\star\star$]\label{exo:m3:on-chain-order-books-and-perpetual-dexs:4}
Un actif a un effet de levier maximal de 20 fois. Quelle est sa marge de maintenance selon l'\cref{dat:m3:on-chain-order-books-and-perpetual-dexs:oracle}, et
à partir de quels fonds propres le coffre reprend-il une position~?
\end{exercise}

\begin{exercise}[$\star\star$]\label{exo:m3:on-chain-order-books-and-perpetual-dexs:5}
Un coffre hérite d'une position vendeuse de 4 millions de dollars à 3 fois le levier maximal, puis le prix double. Quelle est sa perte~?
\end{exercise}

\begin{exercise}[$\star\star$]\label{exo:m3:on-chain-order-books-and-perpetual-dexs:6}
Comparer les risques qu'un fournisseur de liquidité supporte dans un pool à prix d'\omterm{def:m3:blockchains-for-traders:contract}{oracle} et dans le coffre d'une plateforme à carnet d'ordres.
\end{exercise}

\begin{exercise}[$\star\star\star$]\label{exo:m3:on-chain-order-books-and-perpetual-dexs:7}
\emph{Programmation.} Avec \texttt{replay}, porter la probabilité que le preneur arrive en premier à 0.9. Comment les deux règles se comparent-elles~?
\end{exercise}

\begin{exercise}[$\star\star\star$]\label{exo:m3:on-chain-order-books-and-perpetual-dexs:8}
\emph{Repérer l'erreur.} «~La plateforme est décentralisée, donc personne ne peut modifier les règles d'un marché pendant que je détiens une position.~»
\end{exercise}

\section{Problème~: Le coffre écrasé}

\begin{problem}[{Problème du week-end --- un marché étroit, une grosse position vendeuse, un filet de sécurité public}]\label{pb:m3:on-chain-order-books-and-perpetual-dexs:1}
Un \omterm{def:m3:blockchains-for-traders:key}{jeton} peu liquide se négocie comme un perpétuel avec un effet de levier maximal de 3 fois. Un trader ouvre une position vendeuse de 10 millions de dollars et la laisse liquider~; le coffre la reprend à deux
tiers de la marge de maintenance. Le prix spot du \omterm{def:m3:blockchains-for-traders:key}{jeton}, sur de petites plateformes externes, est ensuite poussé à la hausse.

\medskip\noindent\textbf{Partie I --- La reprise.}
\begin{enumerate}
\item Quelle est la marge de maintenance, et quelle marge accompagne la position lorsque le coffre la reprend~?
\item Pourquoi la position ne pouvait-elle pas être clôturée dans le carnet avant la reprise~?
\item Quelle perte le coffre subit-il si le prix monte de 100~\% avant qu'il puisse clôturer~? Et de 400~\%~?
\item Pourquoi le trader peut-il réaliser un profit bien que sa position vendeuse ait été liquidée~?
\item Qu'est-ce que la construction de l'\omterm{def:m3:blockchains-for-traders:contract}{oracle} oblige l'attaquant à faire~?
\end{enumerate}

\medskip\noindent\textbf{Partie II --- Le plafond.}
\begin{enumerate}[resume]
\item Quel plafond sur une position unique maintient la perte du coffre sous 5 millions de dollars pour des compressions allant jusqu'à 500~\%~?
\item À combien la perte aurait-elle été limitée avec un plafond de 2 millions de dollars, à 400~\%~?
\item Pourquoi le plafond doit-il dépendre de la profondeur des marchés spot externes~?
\item Quels autres paramètres la plateforme pourrait-elle resserrer~?
\item Que doit savoir le déposant du coffre sur ces paramètres~?
\end{enumerate}

\medskip\noindent\textbf{Partie III --- Gouvernance.}
\begin{enumerate}[resume]
\item Que fait un vote des \omterm{def:m3:blockchains-for-traders:pow}{validateurs} de radiation aux positions ouvertes~?
\item Qui gagne et qui perd lorsque le règlement-livraison se fait sur la moyenne horaire de l'\omterm{def:m3:blockchains-for-traders:contract}{oracle}~?
\item Le vote est-il un échec de la décentralisation ou l'une de ses caractéristiques~?
\item Comment une \omterm{def:m3:centralised-exchanges:cex}{plateforme centralisée} aurait-elle géré la même attaque~?
\item Comment une société de trading évalue-t-elle ce risque de gouvernance~?
\end{enumerate}

\medskip\noindent\textbf{Partie IV --- Jugement.}
\begin{enumerate}[resume]
\item Un coffre doit-il jamais servir de filet de sécurité pour les marchés peu liquides~?
\item Qu'est-ce qui rend une plateforme on-chain plus sûre qu'une \omterm{def:m3:centralised-exchanges:cex}{plateforme centralisée}, et qu'est-ce qui ne le rend pas~?
\item Où souhaiteriez-vous que la société place son activité de teneur de marché~: sur le DEX à carnet d'ordres, sur le carnet hors chaîne ou dans le pool à prix d'\omterm{def:m3:blockchains-for-traders:contract}{oracle}~?
\item Énoncer le \emph{résultat nommé}~: la perte du coffre en fonction de la compression pour une position de 10 millions de dollars, et le plafond qui la borne à 5 millions de dollars à 500~\%.
\item En une phrase~: qu'est-ce qu'un filet de sécurité public~?
\end{enumerate}
\end{problem}

\section{Questions d'entretien}

\begin{interviewq}[$\star$ \iqroles{trader}]\label{iq:m3:on-chain-order-books-and-perpetual-dexs:1}
Comment une \omterm{def:m3:on-chain-order-books-and-perpetual-dexs:priority}{règle de priorité dans le bloc} qui place les annulations avant les ordres modifie-t-elle l'économie d'un teneur de marché~?
\end{interviewq}

\begin{interviewq}[$\star$ \iqroles{développeur}]\label{iq:m3:on-chain-order-books-and-perpetual-dexs:2}
Qu'est-ce qui change dans votre système de gestion des ordres lorsque la plateforme est une \omterm{def:m3:blockchains-for-traders:chain}{blockchain} à blocs inférieurs à la seconde~?
\end{interviewq}

\begin{interviewq}[$\star\star$ \iqroles{chercheur}]\label{iq:m3:on-chain-order-books-and-perpetual-dexs:3}
Comparez une médiane pondérée des prix externes à une moyenne pondérée par les volumes comme \omterm{def:m3:blockchains-for-traders:contract}{oracle}.
\end{interviewq}

\begin{interviewq}[$\star\star$ \iqroles{risque}]\label{iq:m3:on-chain-order-books-and-perpetual-dexs:4}
Quelles limites fixeriez-vous à l'exposition d'une société au coffre ou au pool d'un \omterm{def:m3:on-chain-order-books-and-perpetual-dexs:perpdex}{DEX de perpétuels}~?
\end{interviewq}

\begin{interviewq}[$\star\star$ \iqroles{trader, chercheur}]\label{iq:m3:on-chain-order-books-and-perpetual-dexs:5}
Comment attaqueriez-vous une \omterm{def:m3:on-chain-order-books-and-perpetual-dexs:oracle}{plateforme à prix d'oracle}, et comment la défendriez-vous~?
\end{interviewq}

\begin{interviewq}[$\star\star\star$ \iqroles{développeur}]\label{iq:m3:on-chain-order-books-and-perpetual-dexs:6}
Concevez un système de rejeu déterministe des carnets d'ordres d'une plateforme à blocs à partir de ses blocs publiés.
\end{interviewq}
